\chapter{Introduction}

Modbench is a Modbus client: it sends requests to Modbus devices (meters,
inverters, drives, PLCs, controllers) and shows what they answer. It reads and
writes registers by hand, or by name with a \emph{device map}, charts values
over time, shows every frame on the wire, and simulates devices when there is
no hardware at hand.

It runs on Windows, Linux and macOS, speaks Modbus over serial lines (RTU and
ASCII) and over the network (TCP, UDP, and RTU frames over TCP or UDP), and
keeps everything on your computer.

\screenshot[7cm]{overview}{The Main page with a device map open.}
  {The whole window: the Main page with a map open on the Device map tab, a test
   running, in the light theme.}

\section{What you can do with it}

\begin{itemize}
  \item \textbf{Send any request}: choose a function code, an address and a
        quantity, press \uiaccent{Start}, and see the values, once or in a loop.
  \item \textbf{Read a device by name}: open a map (\file{.mbmap.json}) and
        every register shows with its name, unit, scaling and state labels.
        Over 400 maps can be downloaded from the shared library.
  \item \textbf{Check that the map fits}: reference values catch a wrong
        address, base or byte order before you trust the other values.
  \item \textbf{Write safely}: typed controls, the map's limits, an optional
        confirmation, and a read-only mode that blocks every write.
  \item \textbf{See what happens}: charts, the raw bytes of each frame
        decoded field by field, a log of every request, and statistics of
        failures by cause.
  \item \textbf{Find and identify devices}: scan unit IDs and serial settings,
        find which addresses answer, ask a device who it is.
  \item \textbf{Simulate devices}: several devices shaped after maps, with
        values that move on their own, to try maps and other programs.
\end{itemize}

\section{A few Modbus words}

If you work with Modbus every day, skip this section.

\begin{description}[style=nextline, leftmargin=0pt, font=\intermedium]
  \item[Client and server] The client (Modbench) asks; the server (the
    device, also called slave) answers. A device never speaks first.
  \item[Unit ID] The address of a device on a serial line, or behind a
    gateway: \val{1} to \val{247}. Over TCP to a single device it is often
    ignored, but must still be sent (usually \val{1} or \val{255}).
  \item[Areas] Each device has four tables of data:
    \begin{center}
    \begin{tabularx}{0.9\linewidth}{l l L}
      \toprule
      \thead{Area} & \thead{Size} & \thead{Read / write} \\
      \midrule
      Coils & 1 bit & read (FC01), write (FC05, FC15) \\
      Discrete inputs & 1 bit & read only (FC02) \\
      Input registers & 16 bits & read only (FC04) \\
      Holding registers & 16 bits & read (FC03), write (FC06, FC16) \\
      \bottomrule
    \end{tabularx}
    \end{center}
  \item[Address] The number of a coil or register in its area, from \val{0}
    to \val{65535} on the wire. Manuals often count from 1, or write the area
    in front (\val{40001} is holding register 1, that is 0 on the wire).
    Modbench shows addresses the way you choose (Section~\ref{sec:notation}).
  \item[Byte and word order] A register holds two bytes; 32- and 64-bit values
    take two or four registers. Devices disagree on which byte and which
    register comes first, so the same bytes can read as very different
    numbers. Maps say which order a device uses; the Data decoder
    (Section~\ref{sec:data-decoder}) helps find it.
  \item[Exception] An answer that says the request failed, with a code:
    \val{02} for an address the device does not have, \val{03} for a value it
    refuses, and so on (Appendix~\ref{app:exceptions}).
\end{description}

\section{The window}
\label{sec:window}

Four pages sit in the title bar. Click one to show it.

\screenshot[3cm]{title-bar}{The title bar and its four pages.}
  {The top of the window: the Modbench logo, the Main, Tools, Library and
   Settings pages, and the window buttons.}

\begin{description}[style=nextline, leftmargin=0pt, font=\intermedium]
  \item[Main] Connect to a device and test it, in three steps: \emph{Connection},
    \emph{Request} and \emph{Test} (Chapter~\ref{ch:main}).
  \item[Tools] The simulator, and tools to identify devices, find addresses,
    decode data and frames, play back recordings, see statistics and set
    alarms (Chapters~\ref{ch:simulator} and~\ref{ch:tools}).
  \item[Library] The device maps on this computer, and the shared online
    library (Chapter~\ref{ch:library}).
  \item[Settings] Theme, text size, safety, what Modbench remembers, updates
    (Chapter~\ref{ch:settings}).
\end{description}

\subsection{Right-click menus}

Most lists and values have a right-click menu with what can be done there:
copy a value or an address, add a register to the chart, write a value, export,
remove. Text fields have \nav{Cut}, \nav{Copy}, \nav{Paste} and
\nav{Select all}.

\subsection{Light and dark}

Modbench follows the system theme, or the one chosen in
\nav{Settings\then Appearance}. The screenshots in this manual use the light
theme.
